\section{Notation and common procedures}\label{sec:notation}

\subsection{Frame, loading and stress quantities}

\block{Frame} A right-handed orthonormal frame at the critical point of a plain specimen:
$x$ is the circumferential (hoop) direction, $y$ the specimen axis, $z$ the outward surface
normal. A thin-walled or solid specimen under axial force, torsion and internal pressure carries
the plane stress $\sigma_{xx}$ (hoop), $\sigma_{yy}$ (axial), $\tau_{xy}$ (torsion); bending is
read as an axial stress $\sigma_{yy}$ at the surface. Stresses are in MPa, nominal and linear
elastic unless a sheet states otherwise.

\block{Harmonic loading} Every test is a constant-amplitude, single-frequency cycle. Each stress
component is written as
\begin{equation}
\Tsig(t) = \Sm + \Ss\sin\omega t + \Sc\cos\omega t ,\qquad 0\le \omega t < 2\pi ,
\label{eq:harmonic}
\end{equation}
with the symmetric tensors $\Sm$ (mean), $\Ss$ and $\Sc$ (in-phase and quadrature parts). For an
axial amplitude $\sigma_a$ and a torsional amplitude $\tau_a$ with phase shift $\delta$,
$\sigma_{yy}(t) = \sigma_m + \sigma_a\sin\omega t$ and
$\tau_{xy}(t) = \tau_m + \tau_a\sin(\omega t + \delta)$, i.e.\ $\tau_a\cos\delta$ enters $\Ss$ and
$\tau_a\sin\delta$ enters $\Sc$. The loading is \emph{proportional} when $\Ss$ and $\Sc$ are
parallel (as six-component vectors); the mean $\Sm$ may then have any direction.

\block{Scalar of one harmonic} A scalar $q(t) = m + a\sin\omega t + b\cos\omega t$ has
\begin{equation}
q_a = \sqrt{a^2+b^2}\ \text{(amplitude, half range)},\qquad q_m = m,\qquad
q_{\max} = m + q_a ,\qquad q_{\min} = m - q_a .
\label{eq:scalar}
\end{equation}
The stress ratio of a uniaxial cycle is $R = q_{\min}/q_{\max}$; fully reversed means $R=-1$.

\block{Plane quantities} A material plane is identified by its unit normal $\n$ ($\n$ and $-\n$
are the same plane). On it, the normal stress and the shear stress vector are
\begin{equation}
\sigma_n(t) = \n^{\mathsf T}\Tsig(t)\,\n ,\qquad
\boldsymbol\tau(t) = \Tsig(t)\,\n - \sigma_n(t)\,\n .
\label{eq:plane}
\end{equation}
$\sigma_{n,a}$, $\sigma_{n,m}$ and $\sigma_{n,\max}$ follow from Eq.~\eqref{eq:scalar}. The shear
amplitude $\tau_a$ of a non-proportional path (the trace of $\boldsymbol\tau(t)$ is an ellipse)
needs a rule, stated by each critical-plane sheet.

\block{Invariants} $I_1 = \tr\Tsig$ (first invariant), $\sigma_H = I_1/3$ (hydrostatic stress),
$\dev\Tsig = \Tsig - \sigma_H\mathbf I$, $J_2 = \tfrac12\,\dev\Tsig:\dev\Tsig$. For one harmonic
the amplitudes are
\begin{equation}
I_{1,a} = \sqrt{(\tr\Ss)^2 + (\tr\Sc)^2},\qquad
\sqrt{J_{2,a}} = \sqrt{\tfrac12\left(\lVert\dev\Ss\rVert^2 + \lVert\dev\Sc\rVert^2\right)} ,
\label{eq:invariants}
\end{equation}
with $\lVert\cdot\rVert$ the Frobenius norm. For one harmonic, the minimum circumscribed ellipse
(Li), the prismatic hull (Mamiya) and the principal-coordinate prismatic hull (Khalij) of the
deviatoric path all give this $\sqrt{J_{2,a}}$; the minimum circumscribed circle does not.
Uniaxial $\sigma_a$: $\sqrt{J_{2,a}}=\sigma_a/\sqrt3$; torsion $\tau_a$: $\sqrt{J_{2,a}}=\tau_a$.

\block{Elastic strains} Where a criterion needs strains, they are the linear-elastic strains of
the stress, $\Teps = \bigl[(1+\nu)\Tsig - \nu\,\tr(\Tsig)\,\mathbf I\bigr]/E$, unless a sheet states
otherwise. Shear strains $\gamma = 2\varepsilon_{ij}$ are engineering strains.

\subsection{Material constants}

\begin{center}\small
\begin{tabular}{@{}lllp{80mm}@{}}
\toprule
Symbol & Code & Unit & Meaning\\\midrule
$\fm$ & \texttt{Sigm1} & MPa & fully reversed ($R=-1$) axial or bending fatigue limit (amplitude)\\
$\tm$ & \texttt{Taum1} & MPa & fully reversed torsional fatigue limit (amplitude)\\
$f_0$ & \texttt{Sig0} & MPa & repeated ($R=0$) axial limit, given as the \emph{maximum} stress of its cycle\\
$t_0$ & \texttt{Tau0} & MPa & repeated ($R=0$) torsional limit, maximum stress\\
$S_u$, $S_y$ & \texttt{Su}, \texttt{Sy} & MPa & ultimate tensile strength, (monotonic) yield strength\\
$E$, $\nu$ & \texttt{E}, \texttt{nu} & MPa, -- & Young's modulus, Poisson's ratio\\
$\kappa$, $r$ & -- & -- & limit ratios $\kappa = \fm/\tm$ and $r = \tm/\fm = 1/\kappa$\\
$N_\text{ref}$ & -- & cycles & life at which the limits are defined ($10^6$--$10^7$)\\
\bottomrule
\end{tabular}
\end{center}
A sheet lists any further constant with its source; a constant fitted by the repository is
stated as such.

\subsection{From a damage parameter to a life}\label{sec:mastercurve}

\block{Damage parameter} Each criterion maps a loading cycle to a scalar damage parameter
$\mathrm{DP}$ (an equivalent stress, a fatigue index, an energy). A fatigue-limit criterion
returns an index that equals 1 on the limit.

\block{Master curve} Unless the method has its own life law, the life follows from
\begin{equation}
\mathrm{DP} = A\,N_f^{\,b},\qquad N_f = (\mathrm{DP}/A)^{1/b},
\label{eq:mastercurve}
\end{equation}
fitted by least squares in $\log$--$\log$ space on the \emph{design set}: the plain, fully reversed
uniaxial tension failures of the material, with $|R+1|\le 10^{-3}$ and every other stress
component below $10^{-3}$ of the peak (a sheet may state the torsion failures instead). Run-outs
never enter the fit. The curve is used only if it is identifiable: at least 5 tests, at least one
decade of life, and a slope significantly different from zero. A criterion that is a fatigue-limit
condition gains finite lives only through Eq.~\eqref{eq:mastercurve}; this is a repository
extension and is stated as such.

\block{Load-equivalent error} The error of a prediction is measured on the load, so that criteria
with different damage parameters are comparable. For a test of life $N_\text{exp}$, $s^*$ is the
factor on all its stresses (means and amplitudes) at which the model predicts $N_\text{exp}$:
$\mathrm{DP}(s^*\Tsig) = A N_\text{exp}^{\,b}$. Then
\begin{equation}
e_L = \ln(1/s^*) ,
\label{eq:eload}
\end{equation}
$e_L>0$ is conservative (the model reaches its curve at a lower load than the test) and the load is
misjudged by $e^{e_L}-1$. A load beyond a criterion's domain counts as infinite damage; a root that
lies on the edge of that domain is rejected.

\block{Fatigue-limit datasets} For a set of limit states (no failures), the reference is the
damage of the fully reversed axial limit, $\mathrm{DP}_\text{ref}=\mathrm{DP}(\fm)$, and
$\mathrm{FI} = \mathrm{DP}/\mathrm{DP}_\text{ref}-1$. The load-scaled index $\mathrm{FI}_\text{load}=1/s-1$,
with $\mathrm{DP}(s\Tsig)=\mathrm{DP}_\text{ref}$, gives the same $e_L=\ln(1+\mathrm{FI}_\text{load})$.

\subsection{Critical-plane search}\label{sec:planes}

A critical-plane criterion evaluates its damage on candidate planes: a 522-plane grid over the
hemisphere of normals, analytical seed planes, and a local refinement. The selected plane maximises
the criterion's own selection metric; ties go to the larger damage. The repository reports the
damage on that plane and, as the variant \texttt{\_ext}, the maximum damage over all planes.

\subsection{Scope}

A criterion is evaluated only on the loading and material its source covers. Outside it the
repository raises a scope error, the model is reported as skipped, and nothing is extrapolated
with an invented formula.
